\chapter{Présentation de la plateforme}
\minitoc

\section{Ce qu'est CIRTDEFENSE}

CIRTDEFENSE est la plateforme d'orchestration autonome de la réponse aux
incidents du Centre de réponse aux incidents informatiques. Elle occupe la
place qui sépare les \emph{capteurs} — pare-feu, sondes réseau, agents de
poste, corrélateur d'événements — des \emph{équipements} sur lesquels on agit
pour contenir une menace.

Sa particularité tient en une phrase : \textbf{elle décide et agit sans
attendre l'accord préalable d'un agent}. L'analyste est informé après coup, et
conserve le pouvoir d'annuler chaque geste appliqué.

\section{Le problème qu'elle traite}

Un centre de réponse national reçoit chaque jour un volume d'alertes qu'aucune
équipe ne peut traiter à la main dans le délai où le confinement a encore un
sens. Le temps qui s'écoule entre la détection d'une exfiltration et le blocage
de l'adresse qui la porte se compte en dizaines de minutes lorsqu'un humain
doit lire, comprendre, décider puis se connecter à chaque équipement. Pendant
ce temps, les données partent.

La réponse habituelle consiste à outiller l'analyste pour qu'il aille plus
vite. CIRTDEFENSE fait un autre choix : \textbf{déplacer la garde humaine de
l'amont vers l'aval}. La plateforme agit, puis rend compte ; l'humain contrôle
au lieu d'autoriser.

\section{Ce qui rend ce choix tenable}

Supprimer le point de contrôle humain avant l'action ne s'improvise pas. Quatre
garde-fous le compensent, et ils sont vérifiés à chaque incident.

\subsection{La réponse doit être documentée}

Avant toute décision, la plateforme cherche dans son fonds documentaire les cas
comparables à ce qu'elle observe. Sans précédent écrit, elle considère qu'elle
agirait sur une hypothèse : \textbf{elle s'abstient et le dit}. C'est le refus
le plus important du système.

\begin{tcolorbox}[colback=green!4!white, colframe=darkgreen,
  title=\textbf{Bon à savoir}, fonttitle=\bfseries, boxrule=0.6pt]
Un taux de refus élevé n'est pas un dysfonctionnement. Il signale que la
documentation de référence est en retard sur les menaces observées. La réponse
est d'enrichir cette documentation, non d'abaisser le niveau d'exigence de la
plateforme.
\end{tcolorbox}

\subsection{Le geste doit être annulable}

Chaque geste que la plateforme sait appliquer est décrit dans un
\emph{catalogue de réversibilité} qui déclare, pour chacun, s'il s'annule
entièrement, partiellement, ou pas du tout, et par quel geste inverse. Un geste
qui se déclare réversible \textbf{doit} porter son verbe d'annulation, sans
quoi la plateforme refuse de le construire.

Seuls les gestes entièrement annulables sont engagés sans humain. Les autres
sont retenus et soumis à décision — c'est l'objet du chapitre~5.

\subsection{L'effet doit être mesuré}

Après un délai de surveillance, la plateforme compare l'état de la machine
avant et après son intervention : joignabilité, latence, taux d'erreur, débit,
sessions actives. Si le confinement a dégradé le service qu'il protégeait,
\textbf{elle l'annule d'elle-même}, sans qu'on le lui demande.

\subsection{Tout doit être rejugeable}

Chaque observation reçue, chaque décision prise, chaque geste appliqué est
inscrit dans un registre dont les entrées sont chaînées par empreintes
cryptographiques. Modifier une entrée passée casse la chaîne à partir de là et
devient visible immédiatement. La base de données refuse d'ailleurs, par
déclencheur, toute mise à jour du journal.

\section{Les profils d'utilisateur}

\begin{table}[h]
\centering
\begin{tabular}{@{}p{2.6cm}p{5cm}p{5.4cm}@{}}
\toprule
\textbf{Profil} & \textbf{Ce qu'il fait} & \textbf{Ce à quoi il accède} \\
\midrule
Analyste &
Suit les interventions, tranche sur les gestes à effet durable, annule ce qui
doit l'être &
Toute l'exploitation, sauf la configuration \\[6pt]
Administrateur &
Installe, configure, rédige les consignes, gère les comptes, réarme le
coupe-circuit &
Tout \\[6pt]
Décideur &
Consulte le portefeuille, édite les rapports &
Lecture, et édition des rapports \\[6pt]
Observateur &
Consulte sans agir &
Lecture seule \\
\bottomrule
\end{tabular}
\caption{Les quatre profils et leur périmètre}
\end{table}

Le premier compte créé à l'installation est un \emph{super-administrateur}~;
il peut ensuite créer les autres comptes. La procédure figure au chapitre~3.

\section{Vocabulaire employé dans ce guide}

Quelques termes reviennent constamment. Ils sont définis ici une fois pour
toutes.

\begin{description}[leftmargin=3.6cm, style=nextline]
  \item[Observation] Ce qu'un capteur signale : une ligne d'alerte, normalisée
        par la plateforme.
  \item[Intervention] Le dossier ouvert par la plateforme pour traiter une
        menace. Il regroupe les observations concordantes et les gestes
        appliqués. Le code l'appelle \emph{incident}.
  \item[Geste] Une action de confinement appliquée à un équipement : bloquer
        une adresse, verrouiller un compte, fermer un port.
  \item[Catalogue métier] Les 22 types d'attaques que le Centre a documentés,
        avec pour chacun sa famille, sa gravité, sa dangerosité et sa priorité.
  \item[Catalogue de réversibilité] Les 38 gestes que la plateforme sait
        appliquer, avec pour chacun son degré d'annulabilité.
  \item[Coupe-circuit] Le dispositif qui suspend toute la réponse automatique
        lorsque la plateforme se trompe trop souvent.
\end{description}
